%! Rational Functions and Vertical Asymptotes — Unit 3, Lesson 3.3 — Group Activity (Answer Key)
\documentclass[10pt]{article}
\usepackage{precalculus-article}
\usepackage{precalculus-key}
\newcommand{\TallMath}[1]{$\displaystyle #1\rule[-1.4em]{0pt}{3.2em}$}

\begin{document}

\pageheader{Unit 3, Lesson 3.3}{Group Activity --- Answer Key}
\namepartnerperiod

\vspace{0.1in}

\begin{scenariobox}[Analyzing Pressure Gauges]{sky}
\small
Engineers monitor pressure in a hydraulic system. The pressure at different points in the
system is modeled by rational functions. A pressure ``spike'' --- where pressure grows without
bound --- corresponds to a vertical asymptote of the model. For each function below, determine
where pressure spikes occur.

\medskip
\renewcommand{\arraystretch}{1.4}
\begin{tabular}{@{} l l @{\qquad} l l @{}}
$r_1(x) = \dfrac{3}{x - 5}$ &
$r_2(x) = \dfrac{x + 2}{x^2 - 4}$ &
$r_3(x) = \dfrac{x^2 - 1}{(x-1)(x+3)}$ &
$r_4(x) = \dfrac{x}{(x+2)^3}$
\end{tabular}
\end{scenariobox}

\vspace{0.1in}

\begin{tcolorbox}[colback=white, colframe=black!40,
    title=\textbf{Tier R --- Remediate},
    fonttitle=\bfseries, arc=2mm, left=3mm, right=3mm, top=2mm, bottom=2mm]

\begin{enumerate}[itemsep=10pt]
  \item $r_1(x) = \dfrac{3}{x - 5}$

        Denominator $= 0$ when $x =$ \ans{$5$}. \quad Vertical asymptote: \ans{$x = 5$}

  \item $r_2(x) = \dfrac{x + 2}{x^2 - 4}$

        Factor the denominator: $x^2 - 4 =$ \ans{$(x-2)(x+2)$}

        Denominator $= 0$ when $x =$ \ans{$2$ and $x=-2$}. \quad
        Vertical asymptote(s): \ans{$x = 2$ and $x = -2$} \textit{(Tier R assumes no common zeros)}

  \item $r_3(x) = \dfrac{x^2 - 1}{(x-1)(x+3)}$

        Denominator $= 0$ when $x =$ \ans{$1$ and $x=-3$}. \quad
        Vertical asymptote(s): \ans{$x = 1$ and $x = -3$} \textit{(Tier R; Tier A/E refine this)}

  \item $r_4(x) = \dfrac{x}{(x+2)^3}$

        Denominator $= 0$ when $x =$ \ans{$-2$}. \quad Vertical asymptote: \ans{$x = -2$}
\end{enumerate}
\end{tcolorbox}

\vspace{0.1in}

\begin{tcolorbox}[colback=white, colframe=black!40,
    title=\textbf{Tier A --- Apply},
    fonttitle=\bfseries, arc=2mm, left=3mm, right=3mm, top=2mm, bottom=2mm]

\begin{enumerate}[itemsep=12pt]
  \item $r_1(x) = \dfrac{3}{x - 5}$: \quad
        Numerator at $x=5$: \ans{$3 \ne 0$}. \quad Result: \ans{VA at $x = 5$}

  \item $r_2(x) = \dfrac{x + 2}{(x-2)(x+2)}$

        At $x = 2$: numerator $=$ \ans{$4 \ne 0$} \quad $\Rightarrow$ \ans{VA at $x = 2$}

        At $x = -2$: numerator $=$ \ans{$0$} \quad $\Rightarrow$ \ans{Potential hole (not a VA); preview for Lesson 1.10}

  \item $r_3(x) = \dfrac{(x-1)(x+1)}{(x-1)(x+3)}$

        Factor the numerator: $x^2 - 1 =$ \ans{$(x-1)(x+1)$}

        At $x = 1$: numerator $=$ \ans{$0$} \quad $\Rightarrow$ \ans{Potential hole (numerator also zero); not a VA}

        At $x = -3$: numerator $=$ \ans{$(-3)^2-1=8\ne0$} \quad $\Rightarrow$ \ans{VA at $x = -3$}

  \item $r_4(x) = \dfrac{x}{(x+2)^3}$: \quad
        At $x = -2$: numerator $=$ \ans{$-2 \ne 0$} \quad $\Rightarrow$ \ans{VA at $x = -2$}
\end{enumerate}
\end{tcolorbox}

\vspace{0.1in}

\begin{tcolorbox}[colback=white, colframe=black!40,
    title=\textbf{Tier E --- Extend},
    fonttitle=\bfseries, arc=2mm, left=3mm, right=3mm, top=2mm, bottom=2mm]

\begin{enumerate}[itemsep=12pt]
  \item \textbf{Explaining a hole.} $r_3(x) = \dfrac{(x-1)(x+1)}{(x-1)(x+3)}$

        \begin{enumerate}[label=(\alph*), itemsep=5pt]
          \item Why is $x = 1$ \emph{not} a vertical asymptote?

                \ans{The factor $(x-1)$ appears in both the numerator and denominator with
                multiplicity 1 in each. Because the numerator's multiplicity (1) is equal to
                the denominator's multiplicity (1) --- not less --- the factor cancels. This
                produces a hole (removable discontinuity) at $x=1$, not a vertical asymptote.}

          \item Confirm that $x = -3$ \textbf{is} a vertical asymptote.

                \ans{At $x=-3$: numerator $= (-3-1)(-3+1) = (-4)(-2) = 8 \ne 0$. The denominator
                is zero but the numerator is not, so $x = -3$ is a vertical asymptote.}
        \end{enumerate}

  \item \textbf{One-sided limits.} $r_4(x) = \dfrac{x}{(x+2)^3}$

        \begin{enumerate}[label=(\alph*), itemsep=6pt]
          \item $x$ slightly greater than $-2$ (e.g., $x = -1.9$):
                Numerator sign: \ans{negative ($-1.9 < 0$)}. \quad
                $(x+2)^3$ sign: \ans{positive (small positive cubed)}. \quad
                So $r_4(x)$ is: \ans{negative}

          \item $x$ slightly less than $-2$ (e.g., $x = -2.1$):
                Numerator sign: \ans{negative ($-2.1 < 0$)}. \quad
                $(x+2)^3$ sign: \ans{negative (small negative cubed)}. \quad
                So $r_4(x)$ is: \ans{positive}

          \item One-sided limits:
                \[
                  \lim_{x \to -2^+} r_4(x) = \ans{$-\infty$}
                  \qquad\qquad
                  \lim_{x \to -2^-} r_4(x) = \ans{$+\infty$}
                \]
        \end{enumerate}
\end{enumerate}
\end{tcolorbox}

\end{document}
